% Internal review appendix only; numeric bands are not experimental results.
\begingroup
\expcolor
\raggedbottom
\section{Internal experiment decision guide}
\label{sec:experiment-decisions}
\textbf{Read red text as work to complete, not a result.}
All result cells remain TBD. The bands below are proposed practical-effect
examples for planning and interpretation. They are not predictions, confidence
intervals, normality limits, or power calculations. Smaller precise effects
can be useful; larger verified effects are not invalid. Actual data determine
both the claim and its scope. Do not stop or select runs because a band is reached.

\paragraph{Completion map.}
E01--E11 retain the earlier handoff IDs. E12--E14 identify the development
behavioral panel, sensitivity, and updating comparisons. E15 adds independent
held-out behavior; E16 audits current applicability. P0 supports the central method
argument; P1 supports scope and efficiency; P2 is needed if claiming the
corresponding extension's empirical benefit. Table~\ref{tab:planned-sequence}
is a work-order guide, not another experiment.

{\small
\begin{tabular}{@{}p{0.13\linewidth}p{0.19\linewidth}p{0.60\linewidth}@{}}
\toprule
ID & Where to fill & Illustrative practical effect, not a forecast \\
\midrule
E01 (P0) & Table~\ref{tab:planned-coverage} & Coverage 80--95\%; valid joint delivery 50--80\% (engineering goals only). \\ [4pt]
E02 (P0) & Table~\ref{tab:planned-content} & $r+W$ vs. $r+H$: Combined +8--17 pp; source observations +4--12 pp. \\ [4pt]
E03 (P0) & Table~\ref{tab:planned-graph} & I--H: +4--12 pp for joint vs. separate checks. C--B: +8--17 pp with the same composer; C--G tests composer choice. \\ [4pt]
E04 (P1) & Table~\ref{tab:planned-utility} & Omitted-requirement utility +8--17 pp; near-zero example $\pm5$ pp. \\ [4pt]
E05 (P1) & Table~\ref{tab:planned-retrieval} & Utility@3 +0.05--0.15; Composable@3 +10--20 pp; Combined +4--12 pp. \\ [4pt]
E06 (P1) & Table~\ref{tab:planned-pool} & Distractor loss at most 5 pp; matched-size source removal is a separate contrast. \\ [4pt]
E07 (P0) & Table~\ref{tab:planned-fourbench} & Graph--Flat +3--8 pp; Graph--No memory +5--12 pp, per benchmark. \\ [4pt]
E08 (P1) & Table~\ref{tab:planned-efficiency} & Time ratio 0.70--0.90; tokens/USD per solve ratio 0.75--0.90 at compatible quality. \\ [4pt]
E09 (P1) & Table~\ref{tab:planned-amortization} & All-in USD/solve ratio 0.75--0.90 at observed valid prefixes; report crossover. \\ [4pt]
E10 (P2) & Table~\ref{tab:planned-agentic} & D--B and E--D +4--12 pp Combined; E--A +8--17 pp as an overall example. \\ [4pt]
E11 (P1) & Table~\ref{tab:planned-method-behavior} & Second-backbone Combined gain +8--17 pp over Flat; evaluate each model separately. \\ [4pt]
E12 (P0) & Table~\ref{tab:planned-method-behavior} & Seven-method Combined gain +8--17 pp; high Preserve alone is not a guarantee. \\ [4pt]
E13 (P1) & Sec.~\ref{sec:planned-coverage} & Vary $k$ and memory budget separately; assess a -5 pp loss margin versus default. \\ [4pt]
E14 (P2) & Sec.~\ref{sec:planned-outcomes} & Updating--frozen +3--8 pp resolution or +4--12 pp Combined; test task-order dependence. \\ [4pt]
E15 (P0) & Table~\ref{tab:planned-heldout} & Independent held-out Combined +5--10 pp; omission burden -2--5 pp with compatible Official. \\ [4pt]
E16 (P0) & Table~\ref{tab:planned-applicability} & False acceptance: proposed 5\% upper-bound tolerance; report useful acceptance and repair effects. \\ [4pt]
\bottomrule
\end{tabular}}

\paragraph{Additional settings without separate result tables.}
Main Table~\ref{tab:planned-main-coverage} summarizes E01/E08/E09; the detailed
E15 diagnostic panel is Table~\ref{tab:planned-heldout-detail}.
E11 repeats all seven E12 arms on the second model. E13 uses five unique
single-factor settings per method (three $k$ values and three token allowances,
sharing the default), comparing CG and Flat with the same total budget.
E14 compares frozen and updating memory within each method under at least
three prespecified task orders, with independent memory stores and permitted
feedback only. These settings require their own result panels before making
the corresponding robustness or adaptation claim.

\clearpage
\subsection{When the observations justify a claim}
\paragraph{Use contrasts, intervals, and compatible quality.}
For a larger-is-better outcome, write $\Delta=\mathrm{CG}-\mathrm{control}$ in
percentage points. A positive estimate with a paired interval clearly above
zero supports an improvement under that protocol. A useful positive estimate
whose interval crosses zero remains uncertain; it does not establish either
superiority or equivalence. For lower-is-better cost use
$\rho=\mathrm{CG}/\mathrm{control}$ and examine whether its interval is below
one while resolution and behavioral quality remain compatible.
Do not infer economic importance from a $p$-value alone.

To claim ``no material degradation,'' choose the tolerance before the new
outcomes. A -5 pp margin is an illustrative design choice here, not a universal
standard: the lower bound of the difference interval must exceed it.
Descriptive Preserve of 95--100\% cannot establish such a bound on its own.
For example, zero failures among 24 independent targets has an exact one-sided
95\% upper failure bound of 11.7\%; repeated attempts on these same targets do
not create 72 independent targets. Boundary cases need appropriate exact or
score intervals, not a degenerate zero-width bootstrap interval.

\paragraph{Respect tasks, repetitions, and missing outcomes.}
For the development panels, report the mean and sample standard deviation of
the three task-level success rates. Estimate the paired uncertainty by
resampling target IDs, keeping all methods and repetitions for each target
together. There are 24 target units, not 72; tests inside a target are
not additional independent observations. One run has no between-run standard
deviation. Do not prescribe or fill an expected standard deviation.
One additional success changes a 24-target repetition by 4.17 pp; the
three-repetition mean changes in steps of 1.39 pp. At the proposed +8--17 pp
scale, the difference is about two to four targets per repetition, so intervals
may still be wide. Set sample size and stopping rules before the runs.

Report resolved/planned, assessed/planned, and jointly assessed pairs.
Use the jointly assessed denominator for paired effects and also disclose
missingness by arm. If $n$ of $N$ planned targets are jointly assessed and their
net paired wins are $d$, a conservative planned-population difference bound is
$[(d-(N-n))/N,\,(d+(N-n))/N]$. Do not silently convert unassessed targets to
failures or remove them from the planned denominator.

\paragraph{The same gain can carry different evidence.}
The following are mathematical examples, not simulated or measured repairs.
With 150 assessed pairs, both 6 wins/0 losses and 10 wins/4 losses give +4.0 pp.
Their two-sided exact McNemar $p$-values are 0.03125 and 0.17957 respectively.
Thus a point estimate inside a proposed band does not guarantee clear evidence.
Prespecify Graph--Flat for the four-benchmark family and apply a correction
such as Holm if making a joint superiority claim; report raw and adjusted
values. A selected best competitor or a list of ablations is not a single
prespecified test. Correct the claimed comparison family or call secondary
comparisons exploratory.

\paragraph{Match the claim to the result.}
A gain over No memory alone supports experience reuse, not graph superiority.
If only F--C improves, the result supports feedback rather than explicit graph
relations. C--B uses the same composer; C--G tests composer choice separately.
The 24-target panels are developmental: E15, with independently authored and
sealed current-contract checks, is needed for the held-out completeness claim.
E03's I--H keeps two check slots and the feedback policy fixed. Its mechanism
claims remain within the development panel, even if E07/E15 also improve.
Group formal outcomes by preregistered adapter support, not successful delivery.
The supported behavioral claim is fewer omissions measured by independent
checks, not satisfaction of every possible requirement. Count preparation
failures and attribute all public-program construction costs to each method.
E16 must test stale requirements, abstention, and the gate's effect on repair.
Better delivery without better Combined supports construction only.
One favorable backbone does not establish model independence. Faster early
resolution with equal final resolution supports speed, not a higher final
success rate. A cost advantage excluding construction supports operating cost,
not amortized cost. Formal cross-benchmark claims require benchmark-specific
results; neither averaging nor source coverage fills a missing benchmark.

\clearpage
\subsection{Existing results that need records, not automatically new runs}
The following are provenance tasks R01--R06 from the earlier handoff. Existing
numeric cells remain black. Recover the original manifests, predictions,
verdicts, configurations, and costs first; a fresh run cannot supply provenance
for an older aggregate.
\begin{itemize}
\item R01: Verify the campaign binding and paired per-task verdicts for
Table~\ref{tab:verified500}; distinguish the current executable method from
historical system versions.
\item R02: Resolve the 35/98 versus 34.69\% discrepancy in
Table~\ref{tab:new-swectx}, and identify the 98/99 task manifest and exclusions.
\item R03: Bind Table~\ref{tab:new-gemini} to its benchmark, sample count,
per-attempt results, and exact pass@k definition.
\item R04: Bind the 236K versus 201K token summary to the same task population
and accounting scope, including failures, retrieval, and construction.
\item R05: Resolve the task-versus-attempt denominators 20 and 30 in
Table~\ref{tab:new-codex55}; compare only matching task sets and definitions.
\item R06: Retain each historical campaign's original denominator, pool,
backbone, budget, and verifier; do not combine them into a new matched panel.
\end{itemize}

\paragraph{Planning rationale and sources.}
The new effect bands are analyst-proposed reference magnitudes. The paper's
existing Verified500 aggregates (+7.4 pp over No memory and +3.8 pp over FAISS)
are a scale reference only; their different protocol does not predict the
new four-benchmark effect. The development bands reflect the coarse resolution
of a 24-target panel. The cost and coverage bands are operational goals rather
than estimates from literature.

Statistical interpretation follows the
\href{https://www.amstat.org/asa/files/pdfs/p-valuestatement.pdf}{ASA statement on $p$-values};
the exact paired examples use the binomial definition documented by
\href{https://www.statsmodels.org/stable/generated/statsmodels.stats.contingency_tables.mcnemar.html}{statsmodels' McNemar reference}.
The importance of interval estimates with few evaluation runs is also discussed
by \href{https://proceedings.neurips.cc/paper/2021/hash/f514cec81cb148559cf475e7426eed5e-Abstract.html}{Agarwal et al. (NeurIPS 2021)}.
None of these sources prescribes the numerical effect bands above.

\paragraph{Reading and exporting this draft.}
The red annotations and this appendix are internal author-review material.
Set \texttt{\string\experimentreviewfalse} in \texttt{experiment-review.tex}
to suppress the extra decision notes and use ordinary text color. This removes
annotations, not the need to complete, qualify, or omit an unfinished empirical
claim. Figure 1 reserves three panels for final-system evidence and remains TBD;
the unchanged historical overview is in Appendix Figure~\ref{fig:historical-overview}.
The two user-supplied design figures and all recorded numerical results are preserved.
\endgroup
